% !PRINTED 90
\weberpage{90}
Wir haben die Konvergenz dieser Reihe für alle Werte von $v$
vorausgesetzt; um diese zu beurteilen, bemerken wir, daß das allgemeine
 Glied dieser Reihe
\[
\begin{array}{lccl}
\text{für} & \nu = +\infty & \text{gleich} & -e^{2\pi i\nu a}\\
\text{\textquotedbl} & \nu = -\infty & \text{\textquotedbl} & e^{-2\pi iv}e^{-2\pi i\nu(\omega-a)}
\end{array}
\]
wird. Es wird also Konvergenz stattfinden, wenn der imaginäre
Teil von $a$ und von $\omega - a$ positiv ist. Setzen wir daher
$\omega = \omega' + i\omega''$, $a = a' + ia''$, so ist die Bedingung der Konvergenz

\weq{8}{0 < a'' < \omega''.}
Wenn man in (7) $a$ durch $a + \frac{1}{2}$ ersetzt, so ergibt sich die
Entwickelung für eine zweite der 16 Funktionen:
\weq{9}{\frac{\vartheta'_{11}\vartheta_{10}(v+a)}{2\pi i\vartheta_{11}(v)\vartheta_{10}(a)} = \sum^{\nu} \frac{(-1)^\nu e^{2\pi ia\nu}}{q^{2\nu}e^{2\pi iv} - 1},}
und wenn man in (7) und (9) $a$ in $a + \frac{1}{2}\omega$ verwandelt:
\weq{10}{\frac{\vartheta'_{11}\vartheta_{01}(v+a)}{2\pi i\vartheta_{11}(v)\vartheta_{01}(a)} = \sum^{\nu} \frac{e^{2\pi ia\nu}q^\nu e^{\pi iv}}{q^{2\nu}e^{2\pi iv} - 1},}
\weq{11}{\frac{\vartheta'_{11}\vartheta_{00}(v+a)}{2\pi i\vartheta_{11}(v)\vartheta_{00}(a)} = \sum^{\nu} \frac{(-1)^\nu e^{2\pi ia\nu}q^\nu e^{\pi iv}}{q^{2\nu}e^{2\pi iv} - 1},}
wobei jedoch zu bemerken ist, daß in den letzten beiden Formeln
die Konvergenzbedingung geändert ist, nämlich:
\[
-\frac{\omega''}{2} < a'' < \frac{\omega''}{2}.
\]
Aus den vier Formeln ergeben sich die übrigen zwölf, wenn
wir $v$ durch $v + \frac{1}{2}$, $v + \frac{1}{2}\omega$, $v + \frac{1}{2}(1+\omega)$ ersetzen, wobei die
Konvergenzbereiche nicht weiter geändert werden. Auf diese Weise
sind die Formeln der Tabelle I am Schlusse des Bandes abgeleitet.

Die so gewonnenen Reihen konvergieren für reelle Werte
von $a$ nicht alle, z. B. tut es nicht die Reihe (7), worin der nach
der Seite der positiven $\nu$ verlaufende Teil
\weq{12}{\sum_{1,\infty}^{\nu} \frac{e^{2\pi i\nu a}}{q^{2\nu}e^{2\pi iv} - 1}}
aufhört zu konvergieren, wenn der imaginäre Teil von $a$ gleich
Null wird, während der andere Teil auch da noch konvergent
%CONT>>
